\section{API Overview}
\label{sec:api-overview}

The S\textsuperscript{4}D system provides a RESTful API built with FastAPI, offering programmatic access to all core functionality. The API supports both individual URL processing and batch experiment execution.

\subsection{Base Configuration}
\label{subsec:base-configuration}

\begin{itemize}
    \item \textbf{Base URL}: \texttt{http://localhost:5000}
    \item \textbf{Content-Type}: \texttt{application/json}
    \item \textbf{Response Format}: JSON
    \item \textbf{CORS}: Enabled for web client integration
    \item \textbf{Documentation}: Available at \texttt{/apidocs} (Swagger UI)
\end{itemize}

\section{Core Processing Endpoints}
\label{sec:core-endpoints}

\subsection{Process URLs}
\label{subsec:process-urls}

\textbf{Endpoint}: \texttt{POST /process-urls}

Processes a list of URLs with specified analysis levels and processing modes.

\subsubsection{Request Format}
\label{subsubsec:request-format}

\begin{verbatim}
{
    "urls": [
        "https://example.com",
        "https://test.com"
    ],
    "levels": [1, 2, 3],
    "mode": 0
}
\end{verbatim}

\textbf{Parameters}:
\begin{itemize}
    \item \texttt{urls} (array): List of URLs to analyze
    \item \texttt{levels} (array): DOM analysis levels (1-10)
    \item \texttt{mode} (integer): Processing mode
        \begin{itemize}
            \item \texttt{0}: Process all URLs together
            \item \texttt{1}: Domain-based processing
        \end{itemize}
\end{itemize}

\subsubsection{Response Format}
\label{subsubsec:response-format}

\begin{verbatim}
{
    "status": "success",
    "message": "URLs processed successfully",
    "clustering_results": [...],
    "processed_clusters": [...],
    "html_files": [...],
    "timing_logs": {...},
    "best_level_info": {...}
}
\end{verbatim}

\textbf{Response Fields}:
\begin{itemize}
    \item \texttt{status}: Success/error indicator
    \item \texttt{message}: Human-readable status message
    \item \texttt{clustering\_results}: HDBSCAN clustering analysis
    \item \texttt{processed\_clusters}: Post-processed cluster data
    \item \texttt{html\_files}: Generated visualization files
    \item \texttt{timing\_logs}: Performance metrics
    \item \texttt{best\_level\_info}: Optimal level identification
\end{itemize}

\subsection{Process Clusters}
\label{subsec:process-clusters}

\textbf{Endpoint}: \texttt{POST /process-clusters}

Processes raw cluster data for visualization and analysis.

\subsubsection{Request Format}
\label{subsubsec:cluster-request}

\begin{verbatim}
{
    "clusters": "[[\"guid1\", \"guid2\"], [\"guid3\", \"guid4\"]]"
}
\end{verbatim}

\textbf{Parameters}:
\begin{itemize}
    \item \texttt{clusters} (string): JSON-encoded cluster arrays
\end{itemize}

\subsubsection{Response Format}
\label{subsubsec:cluster-response}

\begin{verbatim}
{
    "status": "success",
    "message": "Clustering completed successfully",
    "data": "...",
    "sites": {...},
    "html_files": [...],
    "timing_logs": {...},
    "site_similarity": {...}
}
\end{verbatim}

\section{Experiment Management Endpoints}
\label{sec:experiment-endpoints}

\subsection{Run Experiment Batch}
\label{subsec:experiment-batch}

\textbf{Endpoint}: \texttt{POST /experiments/run-batch}

Executes batch experiments with multiple test configurations.

\subsubsection{Request Format}
\label{subsubsec:experiment-request}

\begin{verbatim}
{
    "experiment_name": "Thesis Experiment 1",
    "selected_test_cases": [
        {
            "id": "test_001_0_L1_U3",
            "name": "Test Case 1",
            "mode": 0,
            "levels": [1, 2],
            "urls": ["https://example.com", "https://test.com"],
            "url_count": 2
        }
    ]
}
\end{verbatim}

\subsubsection{Response Format}
\label{subsubsec:experiment-response}

\begin{verbatim}
{
    "status": "success",
    "experiment_summary": {
        "experiment_info": {
            "name": "Thesis Experiment 1",
            "total_tests": 1,
            "successful_tests": 1,
            "failed_tests": 0,
            "success_rate": 100.0,
            "duration": 45.2
        },
        "results": [...],
        "gnuplot_files": [...],
        "generated_images": [...]
    }
}
\end{verbatim}

\subsection{List Experiment Results}
\label{subsec:list-experiments}

\textbf{Endpoint}: \texttt{GET /experiments/results/list}

Retrieves list of completed experiments.

\subsubsection{Response Format}
\label{subsubsec:experiments-list}

\begin{verbatim}
{
    "status": "success",
    "experiments": [
        {
            "id": "20241201_143022_thesis_experiment_1",
            "name": "Thesis Experiment 1",
            "timestamp": "20241201_143022",
            "total_tests": 5,
            "successful_tests": 5,
            "success_rate": 100.0,
            "duration": 120.5
        }
    ],
    "total_count": 1
}
\end{verbatim}

\subsection{Get Experiment Details}
\label{subsec:experiment-details}

\textbf{Endpoint}: \texttt{GET /experiments/results/\{experiment\_id\}}

Retrieves detailed results for a specific experiment.

\subsubsection{Response Format}
\label{subsubsec:experiment-detail}

\begin{verbatim}
{
    "status": "success",
    "experiment_data": {
        "id": "20241201_143022_thesis_experiment_1",
        "experiment_info": {...},
        "results": [...],
        "gnuplot_files": [...],
        "generated_images": [...]
    }
}
\end{verbatim}

\section{File Access Endpoints}
\label{sec:file-endpoints}

\subsection{HTML File Serving}
\label{subsec:html-serving}

\textbf{Endpoint}: \texttt{GET /output\_html\_files/\{filename\}}

Serves generated HTML visualization files.

\textbf{File Structure}:
\begin{itemize}
    \item \texttt{level\_\{level\}/}: Analysis level directories
    \item \texttt{level\_\{level\}\_\{domain\}\_\{identifier\}\_\{guid\}.html}: File naming convention
\end{itemize}

\subsection{Experiment File Downloads}
\label{subsec:experiment-files}

\textbf{Endpoint}: \texttt{GET /experiments/results/\{experiment\_id\}/files/\{filename\}}

Downloads experiment-generated files (GNUplot data, scripts, etc.).

\subsection{Generated Images}
\label{subsec:generated-images}

\textbf{Endpoint}: \texttt{GET /experiments/results/\{experiment\_id\}/images/\{filename\}}

Serves generated chart images (PNG, SVG, etc.).

\section{Error Handling}
\label{sec:error-handling}

\subsection{Standard Error Response}
\label{subsec:error-response}

\begin{verbatim}
{
    "status": "error",
    "message": "Error description",
    "stack_trace": "Detailed error information"
}
\end{verbatim}

\subsection{HTTP Status Codes}
\label{subsec:status-codes}

\begin{itemize}
    \item \texttt{200 OK}: Successful operation
    \item \texttt{400 Bad Request}: Invalid request parameters
    \item \texttt{404 Not Found}: Resource not found
    \item \texttt{500 Internal Server Error}: Server processing error
\end{itemize}

\subsection{Common Error Messages}
\label{subsec:error-messages}

\begin{itemize}
    \item \texttt{"No valid URLs found"}: All provided URLs invalid
    \item \texttt{"ChromeDriver setup failed"}: Browser automation error
    \item \texttt{"Processing timeout"}: Operation exceeded time limit
    \item \texttt{"Cache system error"}: Caching mechanism failure
\end{itemize}

\section{Rate Limiting and Performance}
\label{sec:rate-limiting}

\subsection{Request Limits}
\label{subsec:request-limits}

\begin{itemize}
    \item \textbf{URLs per request}: Maximum 50 URLs
    \item \textbf{Levels per request}: Maximum 10 levels
    \item \textbf{Concurrent requests}: 5 simultaneous requests
    \item \textbf{Timeout}: 30 minutes per request
\end{itemize}

\subsection{Performance Considerations}
\label{subsec:performance-considerations}

\begin{itemize}
    \item \textbf{Processing time}: Scales with URL count × level count
    \item \textbf{Memory usage}: Increases with DOM complexity
    \item \textbf{Cache benefits}: 50-90\% time reduction for cached URLs
    \item \textbf{Optimal batch size}: 5-10 URLs per request
\end{itemize}

\section{Authentication and Security}
\label{sec:authentication}

\subsection{Current Implementation}
\label{subsec:current-auth}

\begin{itemize}
    \item \textbf{Authentication}: None (development/research environment)
    \item \textbf{CORS}: Configured for specific origins
    \item \textbf{Rate limiting}: Built-in request queuing
    \item \textbf{Input validation}: URL format verification
\end{itemize}

\subsection{Production Considerations}
\label{subsec:production-security}

For production deployment, consider implementing:

\begin{itemize}
    \item API key authentication
    \item Request rate limiting per client
    \item Input sanitization and validation
    \item HTTPS encryption
    \item Access logging and monitoring
\end{itemize}

\section{Integration Examples}
\label{sec:integration-examples}

\subsection{Python Client Example}
\label{subsec:python-client}

\begin{verbatim}
import requests
import json

# Process URLs
url = "http://localhost:5000/process-urls"
data = {
    "urls": ["https://example.com", "https://test.com"],
    "levels": [1, 2, 3],
    "mode": 0
}

response = requests.post(url, json=data)
result = response.json()

if result["status"] == "success":
    print(f"Processed {len(data['urls'])} URLs")
    print(f"Generated {len(result['html_files'])} visualizations")
else:
    print(f"Error: {result['message']}")
\end{verbatim}

\subsection{JavaScript Client Example}
\label{subsec:javascript-client}

\begin{verbatim}
async function processUrls(urls, levels, mode = 0) {
    const response = await fetch('http://localhost:5000/process-urls', {
        method: 'POST',
        headers: {
            'Content-Type': 'application/json',
        },
        body: JSON.stringify({
            urls: urls,
            levels: levels,
            mode: mode
        })
    });
    
    const result = await response.json();
    
    if (result.status === 'success') {
        console.log('Processing completed successfully');
        return result;
    } else {
        throw new Error(result.message);
    }
}

// Usage
processUrls(['https://example.com'], [1, 2])
    .then(result => console.log(result))
    .catch(error => console.error(error));
\end{verbatim}

\subsection{cURL Examples}
\label{subsec:curl-examples}

\textbf{Basic URL Processing}:
\begin{verbatim}
curl -X POST "http://localhost:5000/process-urls" \
     -H "Content-Type: application/json" \
     -d '{
       "urls": ["https://example.com"],
       "levels": [1, 2],
       "mode": 0
     }'
\end{verbatim}

\textbf{List Experiments}:
\begin{verbatim}
curl -X GET "http://localhost:5000/experiments/results/list"
\end{verbatim}

\section{API Versioning and Updates}
\label{sec:api-versioning}

\subsection{Current Version}
\label{subsec:current-version}

\begin{itemize}
    \item \textbf{Version}: v1.0
    \item \textbf{Compatibility}: Backward compatible within major version
    \item \textbf{Updates}: Non-breaking changes for minor versions
    \item \textbf{Documentation}: Always reflects current implementation
\end{itemize}

\subsection{Planned Enhancements}
\label{subsec:planned-enhancements}

Future API improvements may include:

\begin{itemize}
    \item WebSocket support for real-time progress updates
    \item Batch processing status endpoints
    \item Advanced filtering and query parameters
    \item Result export in multiple formats (CSV, XML)
    \item Clustering algorithm parameter customization
\end{itemize} 